\documentclass[11pt,a4paper]{article}
\usepackage{amsmath,amssymb,amsthm}
\usepackage[margin=2.5cm]{geometry}
\usepackage{hyperref}

\newtheorem{definition}{Definition}[section]
\newtheorem{theorem}[definition]{Theorem}
\newtheorem{proposition}[definition]{Proposition}
\newtheorem{remark}[definition]{Remark}
\newtheorem{conjecture}[definition]{Conjecture}

\title{Hyperseed Formalization:\\
  The Mind-Blowing/Growing Paraconsistency}
\author{ProtomegaTron (automated formalization)}
\date{2026-09-18}

\begin{document}
\maketitle

\begin{abstract}
We formalize Ben Goertzel's ``The Mind-Blowing/Growing Paraconsistency''
(2021-09-23) within the Hyperseed ontology. Key mappings: paraconsistent
fiber (holding contradictory perspectives without collapse), beauty as
fiber-base resonance, creativity as fiber tension, and PLN as naturally
paraconsistent inference.
\end{abstract}

\tableofcontents

\section{Paraconsistent fiber: holding contradictions}

\begin{theorem}[Fiber must be paraconsistent --- claims 5--8, 16--19, 23]
The fiber bundle's fiber should be paraconsistent --- able to hold
contradictory perspectives without collapsing into triviality:
\[
  F_{\text{paracons}} \models (\phi \wedge \neg\phi)
  \quad \text{without} \quad F_{\text{paracons}} \models \psi \;\forall\psi
\]
Classical fiber would explode from contradiction (ex contradictione quodlibet).
Paraconsistent fiber holds tension productively. This is essential for:
\begin{itemize}
\item Real-world reasoning (incomplete, conflicting, evolving information)
\item Commonsense reasoning (contradictory rules coexisting)
\item Value reasoning (genuinely conflicting values)
\item Multi-perspective reasoning (different perspectives yielding
  contradictory conclusions about the same base)
\end{itemize}
\end{theorem}

\section{Beauty: fiber-base resonance}

\begin{proposition}[Beauty = resonance --- claims 1--4, 24]
Beauty is resonance between fiber and base --- when the fiber's structure
fits the base in a way that transcends any single consistent perspective:
\[
  \text{Beauty}(b) = \text{resonance}(F(b), b)
  \quad \text{requiring paraconsistent } F
\]
Beauty is paradoxical (objective/subjective, evolutionary/transcendent)
precisely because it lives in the paraconsistent zone --- any single
consistent perspective captures only part of it. Full appreciation of
beauty requires paraconsistent fiber that can hold all aspects
simultaneously.
\end{proposition}

\section{Creativity: fiber tension}

\begin{theorem}[Creativity from contradiction --- claims 20--22, 25]
Novel fiber structure emerges from tension between contradictory fiber
sections:
\[
  F_{\text{new}} = \text{resolve}(F_{\phi}, F_{\neg\phi})
  \quad \text{where} \quad F_{\text{new}} \neq F_{\phi}, F_{\neg\phi}
\]
This is Koestler's bisociation formalized: holding two incompatible
fiber sections simultaneously until a new section emerges that
transcends both. Creativity is the fiber's response to its own internal
contradictions --- not just tolerating them but actively exploiting the
tension as fuel for novel structure.
\end{theorem}

\section{Mind and reality: paraconsistency all the way down}

\begin{conjecture}[Ontological paraconsistency --- claims 9--15, 26]
If reality is paraconsistent at fundamental levels (quantum mechanics,
consciousness), then the base space itself is paraconsistent:
\[
  B_{\text{reality}} \text{ is paraconsistent}
  \implies F \text{ must be paraconsistent to faithfully represent } B
\]
Human cognition naturally operates paraconsistently --- we hold contradictory
beliefs productively, experience contradictory emotions simultaneously,
and use contradiction as a driver of exploration and deeper understanding.
This is not a bug but a feature: paraconsistent cognition is more
flexible and creative than classically consistent cognition.
\end{conjecture}

\section{Cross-references}

\begin{remark}[Connections]
\begin{itemize}
\item \textbf{Paraconsistent Interzones} (2021-08-13): paraconsistent
  interzones as spaces where contradictions flourish productively.
\item \textbf{Symbolic Ruminations} (2022-08-06): symbolic/non-symbolic
  paradox requires paraconsistent fiber.
\item \textbf{Parafinity} (2022-01-15): paraconsistency meets parafinity
  at the boundary of classical distinctions.
\item \textbf{Evolving Deeply Ethical AGI} (2026-01-06): paraconsistent
  fiber algebra for value reasoning.
\item \textbf{Evidence Is to Logic What Energy Is to Physics} (2026-03-10):
  evidence conservation as paraconsistent-compatible framework.
\end{itemize}
\end{remark}

\end{document}
